\setcounter{section}{14}

  \zwischenueberschrift{Übungsaufgaben}

\inputaufgabe
{}
{

Es sei $M$ eine
\definitionsverweis {differenzierbare Mannigfaltigkeit}{}{}
mit dem
\definitionsverweis {Kotangentialbündel}{}{}

\mathl{T^*M}{.} Zeige, dass man auf
\mathl{\bigwedge^k T^*M}{} für jedes $k$ eine
\definitionsverweis {Topologie}{}{}
erklären kann, bei der für jede Karte
\maabb {\alpha} { U } { V
} {}
die Abbildung
\maabbdisp {} {\bigwedge^k T^*U} { \bigwedge^k T^*V \cong V \times \bigwedge^k \R^{n*}
} {}
eine
\definitionsverweis {Homöomorphie}{}{}
ist.

}
{} {}
Damit kann man von stetigen und auch von messbaren Differentialformen sprechen.

\inputaufgabe
{}
{

Es sei $M$ eine
$C^2$-\definitionsverweis {differenzierbare Mannigfaltigkeit}{}{}
mit dem
\definitionsverweis {Kotangentialbündel}{}{}

\mathl{T^*M}{.} Zeige, dass
\mathl{\bigwedge^k T^*M}{} für jedes $k$ eine differenzierbare Mannigfaltigkeit ist.

}
{} {}
Damit kann man auch von stetig differenzierbaren Differentialformen sprechen. Allgemeiner gilt: Wenn $M$ eine
$C^m$-\definitionsverweis {Mannigfaltigkeit}{}{}
und

\mavergleichskette
{\vergleichskette
{  \ell
}
{ < }{ m
}
{  }{
}
{    }{
}
{   }{
}
}
{}{}{}
ist, so bezeichnet man mit
\mathl{{ \mathcal E }^{ k }_{\ell } ( U  )}{} die Menge der $\ell$-fach stetig differenzierbaren $k$-Differentialformen.

\inputaufgabe
{}
{

Es sei $M$ eine
\definitionsverweis {differenzierbare Mannigfaltigkeit}{}{}
und
\mathl{{ \mathcal E }^{ k } (  M   )}{} die Menge der
$k$-\definitionsverweis {Formen}{}{}
auf $M$. Zeige, dass ${ \mathcal E }^{ k } (  M   )$ ein
$R$-\definitionsverweis {Modul}{}{}
zu

\mavergleichskette
{\vergleichskette
{R
}
{ = }{ C^0(M,\R)
}
{  }{
}
{    }{
}
{   }{
}
}
{}{}{}
ist.

}
{} {}

\inputaufgabe
{}
{

Es sei $M$ eine
\definitionsverweis {differenzierbare Mannigfaltigkeit}{}{,}

\mavergleichskette
{\vergleichskette
{ P
}
{ \in }{ M
}
{  }{
}
{    }{
}
{   }{
}
}
{}{}{}
ein Punkt und

\mavergleichskette
{\vergleichskette
{ f
}
{ \in }{ C^1(M,\R)
}
{  }{
}
{    }{
}
{   }{
}
}
{}{}{}
eine
\definitionsverweis {stetig differenzierbare Funktion}{}{.}
Es sei

\mavergleichskette
{\vergleichskette
{ v
}
{ \in }{ T_PM
}
{  }{
}
{    }{
}
{   }{
}
}
{}{}{}
ein
\definitionsverweis {Tangentialvektor}{}{,}
der durch einen differenzierbaren Weg
\maabbdisp {\gamma} {{]{-\delta},\delta [}} {M
} {}
mit

\mavergleichskette
{\vergleichskette
{ \gamma(0)
}
{ = }{ P
}
{  }{
}
{    }{
}
{   }{
}
}
{}{}{}
repräsentiert werde. Zeige die Gleichheit

\mavergleichskettedisp
{\vergleichskette
{(df)(P,v)
}
{ =} { (f \circ \gamma)'(0)
}
{ } {
}
{ } {
}
{ } {
}
}
{}{}{.}

}
{} {}

\inputaufgabegibtloesung
{}
{

Es sei $M$ eine
\definitionsverweis {differenzierbare Mannigfaltigkeit}{}{}
und
\maabbdisp {f} { M } {\R
} {}
eine
\definitionsverweis {stetig differenzierbare Funktion}{}{.}
Die Funktion habe im Punkt

\mavergleichskette
{\vergleichskette
{ P
}
{ \in }{ M
}
{  }{
}
{    }{
}
{   }{
}
}
{}{}{}
ein
\definitionsverweis {lokales Extremum}{}{.}
Zeige

\mavergleichskettedisp
{\vergleichskette
{ (df)_P
}
{ =} { 0
}
{ } {
}
{ } {
}
{ } {
}
}
{}{}{.}

}
{} {}

\inputaufgabegibtloesung
{}
{

Es sei $M$ eine
\definitionsverweis {kompakte}{}{}
\definitionsverweis {differenzierbare Mannigfaltigkeit}{}{}
und
\maabbdisp {f} { M } {\R
} {}
eine
\definitionsverweis {stetig differenzierbare Funktion}{}{.}
Zeige, dass es zumindest zwei Punkte auf $M$
\zusatzklammer {vorausgesetzt, dass $M$ zumindest zwei Punkte besitzt} {} {}
gibt, wo die
$1$-\definitionsverweis {Differentialform}{}{}
$df$ verschwindet.

}
{} {}

\inputaufgabe
{}
{

Es sei
\mathl{i: M\subseteq \R^n}{} eine
\definitionsverweis {abgeschlossene Untermannigfaltigkeit}{}{.}
Zeige, dass für eine
\definitionsverweis {differenzierbare Funktion}{}{}
\maabbdisp {f} {\R^n } {\R
} {}
die Beziehung

\mavergleichskettedisp
{\vergleichskette
{ i^* (df)
}
{ =} { d(f \circ i )
}
{ } {
}
{ } {
}
{ } {
}
}
{}{}{}
gilt.

}
{} {}

\inputaufgabe
{}
{

Es sei
\maabbeledisp {f} {\R} {\R
} { t } {f(t)
} {,}
eine
\definitionsverweis {stetig differenzierbare}{}{}
\definitionsverweis {Funktion}{}{}
und es sei

\mavergleichskette
{\vergleichskette
{ \omega
}
{ = }{ g(s)ds
}
{  }{
}
{    }{
}
{   }{
}
}
{}{}{}
eine
$1$-\definitionsverweis {Differentialform}{}{}
auf $\R$. Bestimme
\mathl{f^* \omega}{.}

}
{} {}

\inputaufgabegibtloesung
{}
{

Wir betrachten die Abbildung
\maabbeledisp {\varphi} {\R^3} {\R^2
} {(x,y,z)} {(xyz^2,xy-z^3)
} {.}
Berechne die Matrix der Abbildung
\maabbdisp {\bigwedge^2 T_P (\varphi)} { \bigwedge^2 T_P \R^3 } { \bigwedge^2 T_{\varphi(P)}  \R^2
} {}
im Punkt

\mavergleichskette
{\vergleichskette
{ P
}
{ = }{ (1,3,5)
}
{  }{
}
{    }{
}
{   }{
}
}
{}{}{}
bezüglich einer geeigneten Basis.

}
{} {}

\inputaufgabe
{}
{

Wir betrachten die
\definitionsverweis {stetig differenzierbare Abbildung}{}{}
\maabbeledisp {\varphi} {\R^2 \setminus \{(0,0)\}} {\R^2 \setminus \{(0,0)\}
} {(u,v)} {(u^2,v^3-u)
} {,}
und die
$2$-\definitionsverweis {Differentialform}{}{}

\mavergleichskettedisp
{\vergleichskette
{ \omega
}
{ =} { { \frac{   1  }{  x^2+y^2 } } dx \wedge dy
}
{ } {
}
{ } {
}
{ } {
}
}
{}{}{.}
Bestimme die
\definitionsverweis {zurückgezogene Differentialform}{}{}

\mathl{\varphi^* \omega}{.}

}
{} {}

\inputaufgabegibtloesung
{}
{

Wir betrachten die
$1$-\definitionsverweis {Differentialform}{}{}

\mavergleichskettedisp
{\vergleichskette
{ \omega
}
{ =} { -vdu+udv
}
{ } {
}
{ } {
}
{ } {
}
}
{}{}{}
auf der
$1$-\definitionsverweis {Sphäre}{}{}

\mavergleichskette
{\vergleichskette
{ S^1
}
{ \subset }{ \R^2
}
{  }{
}
{    }{
}
{   }{
}
}
{}{}{,}
wobei die Koordinaten des umgebenden $\R^2$ mit
\mathkor {} {u} {und} {v} {}
bezeichnet seien. Bestimme
\mathl{\pi^* \omega}{} unter der
\definitionsverweis {stetig differenzierbaren}{}{}
Abbildung
\maabbeledisp {\pi} {\R^2 \setminus \{0\}} { S^{1}
} {(x , y) } { { \frac{   1  }{  \sqrt{x^2 + y^2}  } } (x , y)
} {.}

}
{} {}

\inputaufgabegibtloesung
{}
{

Berechne die zurückgezogene Differentialform
\mathl{\varphi^* \tau}{} zu

\mavergleichskettedisp
{\vergleichskette
{ \tau
}
{ =} { dx \wedge dy \wedge dz - w dx \wedge dy \wedge dw + \cos (xy)  dx \wedge dz \wedge dw-yw dy \wedge dz \wedge dw
}
{ } {
}
{ } {
}
{ } {
}
}
{}{}{}
unter der Abbildung
\maabbeledisp {\varphi} {\R^3} {\R^4
} {(r,s,t)} { (r^2s,t, \sin  r ,e^{st}) = (x,y,z,w)
} {.}

}
{} {}

\inputaufgabegibtloesung
{}
{

Es sei
\maabbdisp {\psi} { L } { M
} {}
eine
\definitionsverweis {differenzierbare Abbildung}{}{}
zwischen
\definitionsverweis {differenzierbaren Mannigfaltigkeiten}{}{,}
es sei
\maabbdisp {f} { M } {\R
} {}
eine Funktion auf $M$ und $\omega$ eine
$k$-\definitionsverweis {Form}{}{}
auf $M$. Zeige

\mavergleichskettedisp
{\vergleichskette
{ \psi^*(f \omega)
}
{ =} { \psi^*(f)  \psi^* \omega
}
{ } {
}
{ } {
}
{ } {
}
}
{}{}{.}

}
{} {}

\inputaufgabegibtloesung
{}
{

Es seien

\mavergleichskette
{\vergleichskette
{ W
}
{ \subseteq }{ \R^m
}
{  }{
}
{    }{
}
{   }{
}
}
{}{}{}
und

\mavergleichskette
{\vergleichskette
{ U
}
{ \subseteq }{ \R^n
}
{  }{
}
{    }{
}
{   }{
}
}
{}{}{}
\definitionsverweis {offene Teilmengen}{}{}
und sei
\maabbdisp {\psi} { W } { U
} {}
eine
\definitionsverweis {stetig differenzierbare}{}{}
\definitionsverweis {Abbildung}{}{.}
Es sei
\maabbdisp {f} { U } {\R
} {}
eine stetig differenzierbare Funktion. Folgere aus der
[[Totale Differenzierbarkeit/K/Kettenregel/Fakt|Kettenregel]],
dass

\mavergleichskettedisp
{\vergleichskette
{ d (\psi^*f)
}
{ =} { \psi^*(df)
}
{ } {
}
{ } {
}
{ } {
}
}
{}{}{}
gilt, wobei $\psi^*$ das
\definitionsverweis {Zurückziehen von Differentialformen}{}{}
bezeichnet.

}
{} {}

  \zwischenueberschrift{Aufgaben zum Abgeben}

\inputaufgabe
{6}
{

Es sei $M$ eine
\definitionsverweis {differenzierbare Mannigfaltigkeit}{}{}
mit dem
\definitionsverweis {Kotangentialbündel}{}{}

\mathl{T^*M}{.} Es sei $\omega$ eine
$k$-\definitionsverweis {Differentialform}{}{,} also eine Abbildung
\maabbdisp {\omega} { M } {\bigwedge^k T^*M
} {}
mit
\mathl{\omega(P) \in \bigwedge^k T^*_P M}{} für alle
\mathl{P \in M}{,} wobei dieses Dachprodukt mit der natürlichen Topologie
\zusatzklammer {siehe
[[Mannigfaltigkeit/Dachprodukte des Kotangentialbündels/Topologie/Aufgabe|Aufgabe 14.1]]} {} {} versehen sei. Zeige, dass die folgenden Aussagen äquivalent sind:
\aufzaehlungdrei{ $\omega$ ist
\definitionsverweis {
stetig}{}{.}
}{Für jede
\definitionsverweis {Karte}{}{}
\maabb {\alpha} { U } { V
} {}
mit
\mathl{V \subseteq \R^n}{} und mit der lokalen Darstellung
\mathl{\alpha_* \omega =  \sum_{J,\, { \# \left(   J  \right) } =k} f_J dx_J}{} sind die Funktionen $f_J$ stetig.
}{Es gibt eine offene Überdeckung
\mathl{M= \bigcup_{i \in I} U_i}{} mit
\definitionsverweis {Kartengebieten}{}{} $U_i$ derart, dass in den lokalen Darstellungen
\mathl{\alpha_{i*} \omega =  \sum_{J,\, { \# \left(   J  \right) } =k} f_{i J} dx_J}{} die Funktionen $f_{i J}$ stetig sind.
}

}
{} {}

\inputaufgabe
{4}
{

Es sei
\maabbdisp {\varphi} { L } { M
} {}
eine
\definitionsverweis {differenzierbare Abbildung}{}{}
zwischen zwei
\definitionsverweis {differenzierbaren Mannigfaltigkeiten}{}{}
\mathkor {} {L} {und} {M} {.}
Es seien
\mathkor {} {\omega \in
{ \mathcal E }^{ k } (  M   )} {und} {\tau \in
{ \mathcal E }^{  \ell } (  M   )} {}
\definitionsverweis {Differentialformen}{}{} auf $M$. Zeige die Gleichung

\mavergleichskettedisp
{\vergleichskette
{ \varphi^* ( \omega \wedge \tau)
}
{ =} {\varphi^* \omega \wedge \varphi^* \tau
}
{ } {
}
{ } {
}
{ } {
}
}
{}{}{.}

}
{} {}

\inputaufgabe
{4}
{

Wir betrachten die
\definitionsverweis {stetig differenzierbare Abbildung}{}{}
\maabbeledisp {\varphi} {\R^3 \setminus \{(0,0,0)\}} {N = \left\{ (x,y,z) \in \R^3  \mid   z \neq 0        \right\}
} {(u,v,w)} {(uvw,u^2-vw^5,u^2+v^2+w^2)
} {,}
und die
$2$-\definitionsverweis {Differentialform}{}{}

\mathdisp {\omega =  z^2dx \wedge dy+{ \frac{   xy }{ z } }dx \wedge dz+(xe^y-z)dy \wedge dz} {  }

auf $N$. Bestimme die
\definitionsverweis {zurückgezogene Differentialform}{}{}

\mathl{\varphi^* \omega}{.}

}
{} {}

\inputaufgabe
{6}
{

Wir betrachten die
\definitionsverweis {Einheitssphäre}{}{}

\mavergleichskette
{\vergleichskette
{ S^2
}
{ \subseteq }{ \R^3
}
{  }{
}
{    }{
}
{   }{
}
}
{}{}{,}
wobei die Koordinaten des $\R^3$ mit $x,y,z$ bezeichnet seien. Für welche Punkte

\mavergleichskette
{\vergleichskette
{ P
}
{ \in }{ S^2
}
{  }{
}
{    }{
}
{   }{
}
}
{}{}{}
bilden die Einschränkungen von
\mathkor {} {dx} {und} {dy} {}
auf $S^2$ eine
\definitionsverweis {Basis}{}{}
des
\definitionsverweis {Kotangentialraums}{}{}

\mathl{T^*_PS^2}{.}

}
{} {}

 [[Kategorie:Latexseite]]
